% ============================================================
%  THE ORTYX PROTOCOL
%  Part II — Tension
%  Chapter 6: When Explanation Becomes Atmosphere
% ============================================================

\chapter{When Explanation Becomes Atmosphere}
\label{ch:explanation-atmosphere}

\chapepigraph{A good explanation specifies what it would take to
be wrong. A good atmosphere makes everything feel consistent.
The difference is not always visible from inside either one.}{Flyxion}

\noindent
Part~I followed the Phoenix corpus on its own terms. The result
was a framework of genuine interest: engineering results that
are real, architectural proposals that are serious, and
philosophical claims that connect to live questions in consciousness
studies, cognitive science, and the philosophy of perception.
The excitement that a careful reader feels at the end of Part~I
is appropriate. The framework has earned it.

Part~II presses on a different question: not whether the framework
is interesting, but whether it is structured in a way that makes
it possible to determine, from the outside, whether it is correct.
This is not a question about the framework's internal coherence.
Internally coherent frameworks that resist external evaluation
are possible and common. It is a question about the framework's
relationship to evidence --- about whether there is any
configuration of the world that the framework would recognize
as evidence against its central claims rather than as evidence
for them, or as phenomena requiring reinterpretation rather
than as phenomena confirming the framework's reach.

The transition from Part~I to Part~II is not a transition from
sympathy to hostility. It is a transition from inhabiting the
framework to examining its structure. The examination begins
not with formal audit operators --- those arrive in Part~III
--- but with a set of observations about the texture of the
framework's explanatory activity: the way it moves between
levels, the way it handles anomalies, and the way its
vocabulary functions across different inferential contexts.

\section{Two Kinds of Explanatory Success}
\label{sec:two-kinds}

There are two distinct ways in which a framework can appear
to explain a wide range of phenomena, and they are easy
to confuse from inside the framework.

The first kind is genuine: the framework makes specific
predictions about a domain, those predictions are tested,
and the domain turns out to be structured as the framework
anticipated. The framework's success in one domain then
licenses a cautious extension to adjacent domains, with
appropriate acknowledgment of the increased inferential
distance from the original grounding. This is the normal
progress of science: a framework that works in one place
is tested in other places, and its scope is determined
empirically rather than by declaration.

The second kind is atmospheric: the framework's vocabulary
is sufficiently general that it can be applied to any domain
without generating specific predictions. The application
feels like explanation because it organizes the domain
within the framework's conceptual structure, but it does not
constrain what the domain looks like from within that
structure. Any configuration of the domain is consistent
with the framework, because the framework's terms are defined
broadly enough to accommodate any configuration.

The difference between these two kinds of explanatory success
is not always visible at the level of the framework's
prose. Both produce sentences that sound like explanations.
Both produce conceptual organization that feels satisfying.
The difference is only visible when one asks: given this
framework, what would the domain have looked like if the
framework's central claim were false?

For the Phoenix corpus, the central claim is that temporal
coherence is the primary organizing principle of meaningful
information. What would audio, memory, cognition, and
consciousness look like if that claim were false? This is
not a rhetorical question. It is a structural one. If the
framework can answer it specifically --- if it can describe
observations that would count as evidence against the temporal
coherence hypothesis --- then it is engaged in the first
kind of explanatory success, whatever the current state
of its evidence. If it cannot, it is engaged in the second.

\section{The Vocabulary of Coherence}
\label{sec:vocabulary-coherence}

A recurring feature of the Phoenix corpus is the deployment
of a vocabulary cluster --- coherence, salience, authenticity,
resonance, structure, meaning --- that travels between technical
and non-technical registers without explicit signal that the
register has changed.

In its technical register, ``coherence'' refers to the Marine
jitter metric: a signal is coherent if its peak timing is
stable relative to its running average. This is specific,
measurable, and well-defined. In its experiential register,
``coherence'' refers to the quality of unified, integrated
conscious experience: an experience is coherent if it hangs
together as a whole rather than fragmenting into disconnected
pieces. This is a phenomenological description, not a
measurement.

These are not the same concept. They might be related concepts
--- the hypothesis that temporal coherence in neural dynamics
underlies the phenomenological unity of experience is a
legitimate scientific hypothesis. But the hypothesis is not
the same as the identity. If the corpus treated them as
related concepts whose connection needs to be demonstrated,
the vocabulary would be doing appropriate work. When the
corpus treats them as the same concept under different
descriptions, the vocabulary is doing something else: it
is sustaining the appearance of a tight connection between
the engineering results and the philosophical claims by
using a single word to cover a range of distinct referents.

This is not a deliberate rhetorical strategy. It is a natural
consequence of developing a framework in which the same
fundamental idea --- temporal coherence as the marker of
meaningful structure --- is applied at multiple levels
simultaneously. When a framework's core concept is genuinely
universal, it should appear in every domain. When it merely
\textit{feels} universal because its definition is flexible
enough to fit every domain, the same observation applies
but the inference is different.

\begin{dangerzone}[Symbolic Continuity Without Constraint]
Tracking a single term --- coherence --- across the Phoenix
corpus reveals the following trajectory: it begins as a
time-domain signal property (jitter-based stability), becomes
a feature of biological auditory processing (temporal fine
structure), becomes a property of memory systems (interference
pattern stability), becomes a property of AI architectures
(salience-gated representation), and becomes a property of
conscious experience (phenomenological unity). At each step,
the term carries forward the authority earned at the previous
step without establishing that the same concept is operative.
This is symbolic continuity: the word persists. Whether
inferential continuity is maintained --- whether what makes
a signal coherent and what makes an experience coherent are
related by more than analogy --- has not been demonstrated.
\end{dangerzone}

\section{The Accommodation Problem}
\label{sec:accommodation}

A related phenomenon appears when examining how the corpus
handles observations that might be thought to challenge its
central claims. Consider two such observations.

First: not all high-quality musical experiences require high
temporal coherence in the Marine sense. Noise music, certain
forms of improvised free jazz, and electroacoustic compositions
that deliberately cultivate incoherence and unpredictability
are experienced by sophisticated listeners as intensely
meaningful, emotionally engaging, and aesthetically rich.
If temporal coherence is the primary marker of meaningful
musical structure, how does the framework account for the
meaningfulness of deliberately incoherent music?

Second: many perceptual experiences of high salience and
emotional intensity occur in response to sudden, unexpected,
temporally unpredictable events --- the sound of breaking
glass, an unexpected modulation in a tonal piece, a sudden
dynamic change. These events are precisely high-jitter in
the Marine sense: their timing violates the running model
of the signal. If high salience corresponds to low jitter,
how does the framework account for the salience of the
unexpected?

These are genuine challenges. A framework that responds to
them by generating specific predictions about the boundary
conditions of its claims --- by specifying, for instance,
that the salience of deviation is a second-order effect
that requires a prior coherent baseline, or that the
meaningfulness of incoherent music operates through
different mechanisms than the meaningfulness of coherent
music --- is engaging with the challenges in a way that
allows the framework to be refined by them.

A framework that responds to them by reinterpreting the
challenges as further confirmations --- by arguing, for
instance, that the salience of breaking glass demonstrates
the system's sensitivity to coherence violations, or that
the meaningfulness of noise music reflects the listener's
search for coherence in a signal that frustrates it ---
is doing something different. It is treating the challenges
as occasions for vocabulary application rather than as
tests of the framework's predictions.

\begin{epistemicstatus}{Reconstructive Conjecture}
The Phoenix corpus does not respond explicitly to either
of the challenges posed above. This is not a critique of
the corpus as presently constituted --- it is a research
programme, and research programmes routinely leave boundary
conditions underspecified. The observation is that the
framework's current formulation does not make clear which
of the two response strategies it would employ, which makes
it difficult to determine from outside whether any
observation would count as evidence against the temporal
coherence hypothesis.
\end{epistemicstatus}

\section{The Salience of the Unexpected}
\label{sec:salience-unexpected}

The second challenge --- the salience of unexpected events
--- deserves more careful attention because it touches on
a potential internal tension in the framework.

The Marine salience metric assigns high salience to low-jitter
signals: signals whose timing is stable and predictable.
But salience, in the ordinary sense of the word, is also
high for the unexpected: the sudden, the novel, the
attention-capturing. Psychological research on attention
is unambiguous on this point. The auditory system is
exquisitely sensitive to change, to deviation, to the
unexpected; these properties drive attentional capture
independently of and often in opposition to the stability
that the Marine metric rewards.

The corpus is aware of this tension, at least implicitly.
In some passages, jitter is treated as inverse salience
in the sense that it marks signal degradation rather than
signal value --- high jitter is bad because it indicates
incoherence. In others, the salience metric appears to
be tracking something more like signal richness or
authenticity --- something that high-coherence signals
have and low-coherence signals lack.

These are two different concepts of salience:

\begin{itemize}
  \item \textit{Attentional salience}: the degree to which
  a stimulus captures attention, which correlates positively
  with novelty, deviation, and unpredictability.
  \item \textit{Structural salience}: the degree to which
  a signal carries organized, meaningful structure, which
  the Marine associates with temporal stability and low jitter.
\end{itemize}

Biological auditory systems track both. The auditory cortex
responds strongly to both sustained periodic tones (high
structural salience in the Marine sense) and to transients,
onsets, and frequency glides (high attentional salience in
the standard sense). The two forms of salience are
complementary rather than identical.

The Marine metric, as specified, tracks structural salience.
This is a legitimate engineering choice with clear applications
in audio quality assessment. The tension arises when the
corpus uses the term ``salience'' in contexts where attentional
salience is the relevant concept --- where what matters is
not whether the signal is well-organized but whether it
commands attention. The two concepts come apart precisely
at the moments that are most phenomenologically interesting:
the onset of a phrase, the sudden dynamic shift, the
unexpected harmonic turn.

\section{What the Framework Would Need to Say}
\label{sec:framework-needs}

None of the observations in this chapter constitute refutations
of the Phoenix corpus. They constitute demands: demands for
specification at points where the framework's current
formulation is underspecified in ways that make it difficult
to determine what it predicts.

The framework would be considerably strengthened by explicitly:

\textit{Distinguishing structural from attentional salience}
and specifying when the Marine metric is intended to track
each. This would immediately clarify the boundary conditions
of the metric and generate testable predictions about
cases where the two forms of salience come apart.

\textit{Specifying what temporal incoherence would look like
in a meaningful signal}. Are there meaningful signals that
are genuinely high-jitter in the Marine sense? If yes,
the coherence hypothesis needs to account for them. If no,
the claim that low jitter is necessary for meaningfulness
needs to be defended against the counter-examples above.

\textit{Distinguishing the uses of coherence across levels}.
If ``coherence'' in the signal domain and ``coherence'' in
the phenomenological domain are related by hypothesis rather
than by definition, the hypothesis needs to be stated in
a form that specifies what empirical results would confirm
or disconfirm the relationship.

\textit{Specifying the falsification conditions for the
temporal coherence hypothesis at each level at which it
is applied}. An architecture-level hypothesis and an
experience-level hypothesis may have different falsification
conditions; these need to be kept separate.

These are not unreasonable demands. They are the normal
requirements for a framework at the stage of moving from
research programme to testable theory. The observation
that the Phoenix corpus has not yet met them is not a
judgment that it cannot. It is a statement of where the
framework currently stands in its development.

\section{The Grammar of Atmospheric Frameworks}
\label{sec:atmospheric-grammar}

Frameworks that achieve broad explanatory coverage through
vocabulary generality rather than specific prediction share
a recognizable grammatical structure. Understanding this
structure is useful because it makes the mechanism visible
rather than just its effects.

The grammar works as follows. A key term --- in this case,
coherence --- is introduced with a specific, technical
definition in a domain where it is measurable and useful.
The term then migrates to adjacent domains, carrying with
it the authority established in the original domain.
In the new domains, the term is applied by noting that
the phenomena of those domains can be described using
the vocabulary of the original domain: memory is a form
of resonance, consciousness is a form of coherence,
meaning is a form of temporal structure. Each application
feels like an extension of the framework because the
vocabulary is continuous. But the inferential content of
the term --- what it constrains, what it predicts, what
would count as its absence --- may not carry across with
the word.

The result is a framework that is internally coherent in
a specific sense: all of its claims use the same vocabulary.
But vocabulary coherence is not inferential coherence.
A collection of claims that all use the word ``coherence''
may or may not be making claims that are genuinely related
to one another by more than the word.

\begin{auditprop}
\label{ap:vocabulary-migration}
A framework exhibits vocabulary migration when a term
introduced with a specific, measurable definition in a
source domain is applied in a target domain where its
definition becomes broader, less measurable, or more
ambiguous, without explicit acknowledgment of the change
in definition. Vocabulary migration is a necessary
precondition for atmospheric explanatory expansion: the
appearance of unified explanatory coverage without
demonstrable inferential continuity.
\end{auditprop}

This is the first formal audit proposition in the monograph,
and it is offered tentatively rather than conclusively.
Its status at this stage of the project is that of a
working hypothesis: a candidate characterization of
something observed in the Phoenix corpus that may prove
to have broader applicability. Whether it does --- whether
vocabulary migration is a general feature of high-coherence
speculative systems --- is one of the questions that
Part~III will examine more formally.

\section{The Framework Is Not the Problem}
\label{sec:framework-not-problem}

Before moving to the next chapter, a clarification is
necessary. The observations in this chapter are not
observations about the Phoenix corpus specifically. They
are observations about a class of explanatory behaviour
that appears wherever a framework with genuine local
successes begins extending to broader domains. This class
includes many frameworks that turned out to be substantially
correct as well as many that did not. The presence of
vocabulary migration, of accommodation strategies, and of
underspecified boundary conditions does not establish that
a framework is wrong. It establishes that it has not yet
provided what is needed to determine whether it is right.

What is needed --- and what Part~III will attempt to provide
systematically --- is an audit procedure that can distinguish
between frameworks that are underspecified at an early stage
of development and frameworks that are structurally resistant
to specification. The former can become testable theories;
the latter cannot, and the inability is built into their
structure rather than merely reflecting their current state
of development.

The Phoenix corpus may be the former or the latter. That
determination requires tools that have not yet been assembled.
The present chapter has identified the phenomena that those
tools will need to address. Chapter~7 turns to the first
of those phenomena in detail: the way that formalism can
perform the function of constraint without actually
providing it.

\medskip

The tension is now accumulating. The reader should feel,
at this point, something like the experience of watching
a tightrope walker move from a stable platform onto the
rope itself: the motion is smooth, the posture is confident,
but the support structure has changed. The framework is
no longer standing on engineering results. It is standing
on vocabulary. Whether vocabulary can bear that weight
is the question Chapter~7 begins to press.
